\section{Variabel Penelitian}

Penetapan variabel penelitian dilakukan agar aktivitas \textit{Demonstration} dan \textit{Evaluation} dapat dioperasionalkan secara jelas serta pengujian hipotesis berlangsung secara terkontrol. 
Variabel penelitian dibagi menjadi variabel bebas, terikat, dan kontrol.

Variabel bebas mencakup tiga faktor. 
Pertama, jenis algoritma yang terdiri atas tujuh kondisi, yaitu CoDA-EDA, \textit{compact-GA}, \textit{compact-PSO}, BOA, iBOA, \textit{Ant Colony Optimization} (ACO), dan \textit{Upgraded Black-winged Kite} (UBK). 
Kedua, tingkat korelasi antaratribut pada dataset pengujian, yang dikendalikan pada kategori rendah, sedang, dan tinggi melalui perluasan generator ABAC Lab. 
Ketiga, skenario \textit{drift} data akses, yaitu besarnya perubahan antarsiklus yang memengaruhi pemicu pembaruan struktur dependensi pada CoDA-EDA.

Variabel terikat mencakup empat kelompok metrik: (1) kualitas kebijakan, meliputi akurasi/\textit{F1-score}, \textit{Weighted Structural Complexity} (WSC), dan \textit{coverage}; (2) jejak memori, yang diukur melalui konsumsi memori puncak; (3) waktu komputasi, mencakup waktu pembelajaran awal dan satu siklus pembaruan inkremental; serta (4) kinerja re-optimasi pada skenario \textit{drift}, yang membandingkan waktu dan kualitas kebijakan antara strategi \textit{warm-start} dan \textit{re-mining} penuh.

Untuk menjaga keadilan perbandingan, seluruh algoritma dijalankan pada perangkat keras yang sama, yaitu Raspberry Pi 5, menggunakan dataset identik untuk setiap skenario. 
Kriteria pemberhentian disetarakan antaralgoritma, sedangkan setiap kombinasi algoritma dan skenario dieksekusi berulang dengan \textit{random seed} yang berbeda.

Definisi operasional seluruh variabel beserta keterkaitannya dengan hipotesis dirangkum pada \ref{tab:definisi-operasional}.

\begin{longtable}{|c|>{\raggedright\arraybackslash}p{3cm}|c|>{\raggedright\arraybackslash}p{5cm}|>{\centering\arraybackslash}p{3cm}|}
\caption{Definisi Operasional Variabel Penelitian}
\label{tab:definisi-operasional} \\
\hline
\textbf{No} & \textbf{Variabel} & \textbf{Jenis} & \textbf{Definisi Operasional} & \textbf{Terkait Hipotesis} \\
\hline
\endfirsthead

\multicolumn{5}{c}{{\bfseries \tablename\ \thetable{} -- lanjutan dari halaman sebelumnya}} \\
\hline
\textbf{No} & \textbf{Variabel} & \textbf{Jenis} & \textbf{Definisi Operasional} & \textbf{Terkait Hipotesis} \\
\hline
\endhead

\hline
\multicolumn{5}{r}{Lanjut ke halaman berikutnya} \\
\endfoot

\hline
\endlastfoot

% ==================== DATA BARIS ====================
1 & Algoritma yang dibandingkan & Bebas & Tujuh kondisi: CoDA-EDA (diusulkan) dan enam algoritma pembanding (\textit{compact-GA, compact-PSO}, BOA, iBOA, ACO, UBK). Dari keenam pembanding, tiga berpopulasi eksplisit (BOA, ACO, UBK), sedangkan \textit{compact-GA/PSO} dan iBOA berjejak memori datar terhadap NP. & H1 (vs \textit{compact-GA/PSO}); H2 (vs BOA, ACO, UBK); H3 (vs iBOA) \\
\hline
2 & Tingkat korelasi atribut & Bebas & Tiga tingkat terkontrol (rendah, sedang, tinggi) yang disuntikkan pada generator ABAC Lab yang diperluas. & H1 \\
\hline
2b & Dimensi ruang solusi (D) & Bebas & Banyaknya pasangan atribut-nilai yang membentuk vektor keputusan: D = 436 pada \textit{workforce} dan D = 1.907 pada \textit{e-document}. Ditetapkan sebagai variabel bebas berdasarkan studi pendahuluan yang menunjukkan dimensi soal sebagai faktor penentu keunggulan pemodelan dependensi. & H1 \\
\hline
3 & Skenario \textit{drift} data akses & Bebas & Besaran perubahan (delta) pada log akses antar-siklus pembaruan yang memicu pembaruan struktur dependensi (tidak ada/rendah/sedang). & H3, H4 \\
\hline
4 & Kualitas kebijakan & Terikat & Akurasi/\textit{F-score}, \textit{Weighted Structural Complexity} (WSC), dan cakupan (\textit{coverage}) aturan terhadap log akses; formula rinci pada bagian metode evaluasi. & H1, H3, H4 \\
\hline
5 & Jejak memori & Terikat & Konsumsi memori puncak (\textit{peak resident set size}) selama algoritma berjalan pada perangkat target; prosedur pengukuran pada bagian metode evaluasi. & H2 \\
\hline
6 & Waktu komputasi & Terikat & Waktu pembelajaran awal (\textit{initial training}) dan waktu satu siklus pembaruan inkremental; prosedur pengukuran pada bagian metode evaluasi. & H3 \\
\hline
7 & Efisiensi re-optimasi pada \textit{drift} & Terikat & Perbandingan waktu dan kualitas kebijakan antara strategi \textit{warm-start} dan \textit{re-mining} penuh dari awal pasca-\textit{drift}. & H4 \\
\hline

\end{longtable}
